%%%%%%%%%%%%%%%%%%%%%%%%%%%%%%%%%%%%%%%%%%%%%%%%%%%%%%%%%%%%%%%%%%%%%%%%%%%%%%%%
%% Replacement LaTeX for the market-condition analysis.
%%
%% Drop Section A into the methodology, replace the old Table 7 block with
%% Section B, and replace the corresponding paragraphs in the results and
%% limitations sections with Sections C and D.
%%
%% Every number quoted here comes from results/summary.json, which is generated
%% by aggregate.py from the per-fold prediction files in results/preds/.
%%%%%%%%%%%%%%%%%%%%%%%%%%%%%%%%%%%%%%%%%%%%%%%%%%%%%%%%%%%%%%%%%%%%%%%%%%%%%%%%


%% ===========================================================================
%% SECTION A - new methodology subsection
%% ===========================================================================
\subsection{Market regimes and evaluation protocol}
\label{sec:regimes}

The comparison in Section~\ref{sec:results} is reported separately for three
market regimes. This subsection states how those regimes are defined and how the
test data are constructed, since both differ from a single pooled train/test
split.

\noindent\textbf{Blocked walk-forward validation.}
Sequences of consecutive trading days overlap heavily: with a look-back of
$L = 20$ days, adjacent input windows share 19 of their 20 days. A randomly
shuffled train/test split therefore places near-duplicate windows on both sides
of the split, which flatters any model that can interpolate between them. We
instead use blocked walk-forward validation with five folds. Let $n$ be the
number of sequences available at horizon $h$. Fold $i$ trains on sequences
$[0, s_i)$ and tests on the contiguous block $[s_i, s_{i+1})$, where
$s_0 = \lfloor 0.4 n \rfloor$ and the remaining sequences are divided into five
equal blocks. No test sequence ever precedes its own training data, and the
pooled test blocks span the last 60 percent of the sample period, so the
evaluation covers several distinct market episodes rather than one.

\noindent\textbf{Regime definition.}
Each test window is labelled from the $L$ observed days the model receives as
input, so no information from the forecast period is used to assign a label.
Writing $c_1, \dots, c_L$ for the scaled closing prices in the window and
$\Delta_t = c_{t+1} - c_t$, we compute

\begin{equation}
v = \operatorname{sd}(\Delta_1, \dots, \Delta_{L-1}),
\qquad
\mathrm{ER} = \frac{\lvert c_L - c_1 \rvert}{\sum_{t=1}^{L-1} \lvert \Delta_t \rvert},
\label{eq:regime_stats}
\end{equation}

where $v$ measures realised volatility over the window and $\mathrm{ER}$ is the
Kaufman efficiency ratio, which is close to 1 when price movement is
directional and close to 0 when it is choppy. Both statistics are invariant to
the affine Min-Max transform applied to the series, so they are unaffected by
the scaling choice discussed in Section~\ref{sec:limitations}.

Thresholds are the 75th percentiles $v^{\ast}$ and $\mathrm{ER}^{\ast}$ of the
two statistics computed on the training portion of each fold only, never on the
test block. Labels are then assigned in this order:

\begin{align}
\text{Volatile} \quad &\text{if } v > v^{\ast},\notag\\
\text{Trending} \quad &\text{else if } \mathrm{ER} > \mathrm{ER}^{\ast},
\label{eq:regime_rule}\\
\text{Normal}   \quad &\text{otherwise.}\notag
\end{align}

The rule is exhaustive and mutually exclusive by construction. The resulting
window counts are reported in Table~\ref{table_regime_comparison}; roughly
55~percent of test windows are Normal, 30~percent Volatile and 15~percent
Trending.

\noindent\textbf{Identical test sets and a naive baseline.}
At every horizon both architectures are trained on the same training blocks and
scored on the same test windows, so the comparison is like-for-like. We also
report a persistence baseline that repeats the last observed closing price for
all $h$ steps ahead. Without such a baseline it is not possible to say whether a
given RMSE reflects genuine forecasting skill or merely the fact that daily
prices change slowly.


%% ===========================================================================
%% SECTION B - replacement for the old Table 7
%% ===========================================================================
%% Paste the contents of results/table_regime_comparison.tex here. It is
%% generated rather than hand-written so that the table cannot drift out of
%% step with the logged runs.
\begin{table}[!ht]
\centering
\scriptsize
\renewcommand{\arraystretch}{1.15}
\setlength{\tabcolsep}{4pt}
\begin{tabular}{|l|l|c|c|c|c|c|}
\hline
\textbf{Horizon} & \textbf{Regime} & \textbf{Windows} & \textbf{LSTM RMSE} & \textbf{KAN RMSE} & \textbf{Naive RMSE} & \textbf{LSTM vs KAN} \\
\hline
  \multirow{3}{*}{1-Day} & Normal & 277 & $0.055 \pm 0.011$ & $0.051 \pm 0.005$ & 0.020 & 0.9$\times$ \\
   & Volatile & 175 & $0.085 \pm 0.014$ & $0.057 \pm 0.010$ & 0.020 & 0.7$\times$ \\
   & Trending & 79 & $0.063 \pm 0.010$ & $0.063 \pm 0.002$ & 0.019 & 1.0$\times$ \\
  & \textit{All} & 531 & $0.068 \pm 0.009$ & $0.055 \pm 0.006$ & 0.020 & 0.8$\times$ \\
\hline
% horizon 1 used LSTM configuration 100u-linear
  \multirow{3}{*}{2-Day} & Normal & 277 & $0.053 \pm 0.002$ & $0.058 \pm 0.004$ & 0.026 & 1.1$\times$ \\
   & Volatile & 176 & $0.096 \pm 0.005$ & $0.069 \pm 0.002$ & 0.024 & 0.7$\times$ \\
   & Trending & 78 & $0.069 \pm 0.004$ & $0.068 \pm 0.003$ & 0.023 & 1.0$\times$ \\
  & \textit{All} & 531 & $0.072 \pm 0.003$ & $0.063 \pm 0.001$ & 0.025 & 0.9$\times$ \\
\hline
% horizon 2 used LSTM configuration 10u-tanh
  \multirow{3}{*}{100-Day} & Normal & 197 & $0.110 \pm 0.005$ & $0.120 \pm 0.006$ & 0.107 & 1.1$\times$ \\
   & Volatile & 170 & $0.105 \pm 0.006$ & $0.122 \pm 0.005$ & 0.118 & 1.2$\times$ \\
   & Trending & 105 & $0.117 \pm 0.019$ & $0.132 \pm 0.008$ & 0.107 & 1.1$\times$ \\
  & \textit{All} & 472 & $0.110 \pm 0.009$ & $0.124 \pm 0.006$ & 0.111 & 1.1$\times$ \\
\hline
% horizon 100 used LSTM configuration 10u-tanh
  \multirow{3}{*}{200-Day} & Normal & 187 & $0.153 \pm 0.011$ & $0.129 \pm 0.002$ & 0.139 & 0.8$\times$ \\
   & Volatile & 144 & $0.113 \pm 0.004$ & $0.113 \pm 0.000$ & 0.132 & 1.0$\times$ \\
   & Trending & 81 & $0.206 \pm 0.014$ & $0.149 \pm 0.003$ & 0.115 & 0.7$\times$ \\
  & \textit{All} & 412 & $0.153 \pm 0.010$ & $0.128 \pm 0.002$ & 0.132 & 0.8$\times$ \\
\hline
% horizon 200 used LSTM configuration 10u-tanh
\end{tabular}
\caption{Regime-conditional performance of LSTM and KAN across forecast horizons (RMSE in Min-Max scaled space, lower is better). Regimes are assigned from the 20-day input window using the rule defined in Section~\ref{sec:regimes}; \emph{Windows} is the number of test windows falling in each regime. Figures are the mean and sample standard deviation over three random seeds, each pooled across five blocked walk-forward folds. \emph{Naive} is the persistence baseline that repeats the last observed close price. Both architectures are evaluated on identical test windows at every horizon, including 200 days.}
\label{table_regime_comparison}
\end{table}



%% ===========================================================================
%% SECTION C - replacement results narrative
%% ===========================================================================
Table~\ref{table_regime_comparison} reports RMSE by regime and horizon under the
protocol of Section~\ref{sec:regimes}. Three observations follow.

\noindent\textbf{The two architectures perform comparably.}
At the 1-day horizon the KAN attains a pooled RMSE of $0.055 \pm 0.006$ against
$0.068 \pm 0.009$ for the LSTM; at 2 days the figures are $0.063 \pm 0.001$ and
$0.072 \pm 0.003$. At 100 days the ordering reverses, with the LSTM at
$0.110 \pm 0.009$ against $0.124 \pm 0.006$, and at 200 days it reverses again.
In every case the gap is within roughly one and a half standard deviations of
the seed-to-seed spread. We therefore do not find evidence that either
architecture dominates the other on this dataset under matched evaluation.

\noindent\textbf{Regime structure is present but modest.}
Both models degrade in the Volatile regime relative to Normal conditions at
short horizons, and the LSTM degrades most: its 2-day RMSE rises from
$0.053 \pm 0.002$ under Normal conditions to $0.096 \pm 0.005$ under Volatile
conditions, an 80~percent increase, while the KAN rises from
$0.058 \pm 0.004$ to $0.069 \pm 0.002$. At the 200-day horizon the pattern
inverts, with Trending windows the hardest for both models. The regime
breakdown is thus informative, but the differences between regimes are smaller
than the differences between horizons.

\noindent\textbf{Neither model beats persistence at short horizons.}
The naive baseline attains RMSE $0.020$ at 1 day and $0.025$ at 2 days,
roughly three times lower than either trained model. At 100 and 200 days the
three methods are comparable, with the baseline still competitive
($0.111$ and $0.132$ respectively). This is the most consequential result in
the study: under matched evaluation, neither architecture as configured here
extracts forecasting skill beyond the persistence of the price level. Any claim
about the relative merits of KAN and LSTM has to be read in that light.


%% ===========================================================================
%% SECTION D - replacement limitations text
%% ===========================================================================
\noindent\textbf{Scaling.} Min-Max scaling is fitted over the full sample
before splitting, so the scaler carries information from the whole period. All
errors are reported in this scaled space, which makes them comparable across
configurations but not interpretable in currency units. Refitting the scaler on
each training block is the correct treatment and is left to future work.

\noindent\textbf{Training budgets remain unmatched.} The KAN is trained with 10
full-batch L-BFGS iterations and the LSTM with 25 epochs of mini-batch RMSprop,
as in the original specification. The two therefore receive different numbers of
effective parameter updates, and the comparison should not be read as a
statement about the architectures under equal compute.

\noindent\textbf{Single instrument.} All results are for one ASX-listed equity
over 2020--2023. The regime rule of Equation~\eqref{eq:regime_rule} uses
relative thresholds and so transfers to other series, but the conclusions are
not established beyond this one.

\noindent\textbf{Test spans differ across horizons.} A target window of $h$ days
consumes the last $h$ observations, so the pooled test period shrinks as the
horizon grows: 2021-06 to 2023-07 at 1 and 2 days, 2021-04 to 2023-03 at 100
days, and 2021-02 to 2022-10 at 200 days. Comparisons between architectures at a
given horizon are like-for-like, but comparisons of the same architecture across
horizons are made on partly different periods.

\noindent\textbf{Model capacity.} Both architectures are small: the LSTM
configurations follow the original grid and the KAN uses the width heuristic of
the original implementation. The finding that neither beats persistence may
reflect these specific configurations rather than the architectures in general.
